% Computational discovery chapter: partitions, Bernoulli numbers, and ODEs
\chapter{Partitions, Power Sums, and Numerical Differential Equations}
\label{ch:discovery_partitions_calculus}

This chapter follows a recurring Eulerian pattern: calculate small cases, look for
structure, and only then formulate a theorem. The computer is an experimental
instrument; the proof remains a mathematical argument.

\section{Integer Partitions}

A partition of $n$ is an unordered sum of positive integers. Let $p(n)$ count
partitions of $n$. Direct enumeration begins:

\[
p(0),p(1),\ldots,p(10)=1,1,2,3,5,7,11,15,22,30,42.
\]

\begin{lstlisting}[language=Python]
from euler import partition_function_dp

p = partition_function_dp(20)
print(list(enumerate(p)))
\end{lstlisting}

The surprising recurrence uses the generalized pentagonal numbers
$g_k=k(3k-1)/2$:

\begin{theorem}[Euler's pentagonal recurrence]
\[
p(n)=\sum_{k\geq1}(-1)^{k-1}\bigl(p(n-g_k)+p(n-g_{-k})\bigr),
\]
where $p(0)=1$ and $p(m)=0$ for $m<0$.
\end{theorem}

The computation suggests the theorem before the proof. Euler's generating-function
identity is

\[
\prod_{m=1}^{\infty}(1-x^m)=1+\sum_{k=1}^{\infty}(-1)^k
\left(x^{k(3k-1)/2}+x^{k(3k+1)/2}\right).
\]

\begin{lstlisting}[language=Python]
from euler import euler_partition_theorem

for n in range(1, 16):
    result = euler_partition_theorem(n)
    print(n, result['q(n)'], result['p_odd(n)'], result['equal'])
\end{lstlisting}

This verifies another Euler discovery: partitions into distinct parts and partitions
into odd parts occur in equal numbers. The natural next question is to find the
bijection, not merely to observe equality.

\section{Bernoulli Numbers and Power Sums}

The coefficients in Euler--Maclaurin summation also organize sums of powers. The
Bernoulli numbers are defined by

\[
\frac{x}{e^x-1}=\sum_{m=0}^{\infty}B_m\frac{x^m}{m!},
\qquad B_0=1,\quad B_1=-\frac12.
\]

\begin{lstlisting}[language=Python]
from euler import bernoulli_numbers, faulhaber_sum

print(bernoulli_numbers(10))
for k in range(1, 5):
    print(k, faulhaber_sum(10, k))
\end{lstlisting}

\begin{theorem}[Faulhaber's formula]
\[
\sum_{r=1}^{n}r^k=\frac{1}{k+1}\sum_{j=0}^{k}(-1)^j
\binom{k+1}{j}B_j n^{k+1-j}.
\]
\end{theorem}

The program uses exact rational arithmetic, so the computer can verify identities
without hiding rounding error. Students should first guess the polynomial for
$1+2+\cdots+n$, then for squares and cubes, before meeting the general formula.

\section{Euler's Method}

For $y'=f(x,y)$, Euler's method replaces the curve locally by its tangent line:

\[
y_{r+1}=y_r+h f(x_r,y_r).
\]

For $y'=y$, $y(0)=1$, the exact answer at $x=1$ is $e$. The following experiment
reveals first-order convergence, while Heun and RK4 reveal higher orders.

\begin{lstlisting}[language=Python]
from euler import euler_method, improved_euler_method, runge_kutta_4
from math import exp

f = lambda x, y: y
for method in (euler_method, improved_euler_method, runge_kutta_4):
    errors = []
    for h in (0.2, 0.1, 0.05):
        value = method(f, 0.0, 1.0, h, round(1/h))[-1][1]
        errors.append(abs(value - exp(1)))
    print(method.__name__, errors, errors[0]/errors[1])
\end{lstlisting}

The ratios approach $2$, $4$, and $16$, respectively. This is a computational route
to the ideas of local error, global error, and convergence order.

\section{Exercises}
\begin{exercise}
Compute $p(n)$ by enumeration for $n\leq 10$ and compare it with Euler's recurrence.
\end{exercise}
\begin{exercise}
Use exact arithmetic to prove the formulas for the first four power sums.
\end{exercise}
\begin{exercise}
Repeat the ODE experiment for $y'=-2y$ and explain why large step sizes become unstable.
\end{exercise}
